\chapter{Text the Harness Sends to the Model}
\label{app:prompts}

Every piece of text the harness itself says to the model is in \code{harness/prompts.py} (or, for the ending tools,
\code{read_result} and \code{run_python}, next to their schemas). The texts below are the current version,
\code{PROMPT_VERSION} = v7, extracted from the code on 9 October 2026; line breaks are added for
printing. The tool descriptions are in Appendix~\ref{app:tools}.

% source: harness/prompts.py, harness/answers.py, harness/tools/handles.py, harness/sandbox/__init__.py

\section*{System prompt (\code{SYSTEM_PROMPT})}
Sent as the first message of every run.
\begin{lstlisting}
You answer questions about a graph G. You cannot see G directly; use the tools to
inspect it. When you know the answer, call submit_answer. If the question cannot be
answered with the available tools, call cannot_answer instead of guessing.
\end{lstlisting}

\section*{Nudge (\code{NUDGE})}
Sent when the model replies with plain text instead of calling a tool.
\begin{lstlisting}
You did not call a tool. Use the tools to inspect G, call submit_answer with your final
answer, or call cannot_answer if the question cannot be answered.
\end{lstlisting}

\section*{Format error (\code{FORMAT_ERROR})}
Sent when the plain text looks like a hand-written tool call that was not executed (in strict mode, or in lenient mode when the text cannot be parsed into a tool call).
\begin{lstlisting}
Your last message contains a tool call written as plain text, so it was NOT executed. Do
not write tool calls in your reply; call the tool through the function-calling
interface.
\end{lstlisting}

\section*{Empty reply (\code{EMPTY_REPLY})}
Sent when the reply has no text and no tool call (from v6).
\begin{lstlisting}
Your reply was empty: no text and no tool call. If you already know the answer, call
submit_answer now. Otherwise make your next tool call, or call cannot_answer if the
question cannot be answered.
\end{lstlisting}

\section*{Code strategy hint (\code{CODE_HINT})}
Added to the system prompt only when \code{run_python} is offered and the hint is switched on (setup H).
\begin{lstlisting}
How to work: use a direct tool for a single lookup. When the answer needs many lookups
(a value for every neighbor, every node, or every pair of nodes) or a comparison over
many results, write ONE run_python call that calls the tool functions in a loop and puts
the answer in `result`: it is faster and avoids mistakes in long lists. You can combine
both: a tool call to find the candidates, then code to check them all.
\end{lstlisting}

\section*{Compaction notes (\code{SHORTENED_RESULT}, \code{SHORTENED_CODE})}
What an older tool result or older code is prefixed with or replaced by when the conversation is compacted (from v7).
\begin{lstlisting}
(shortened; call the tool again for the full result)
# (older code shortened to save space)
\end{lstlisting}

\section*{\code{submit_answer} (answer type \code{yes_no_with_cycle})}
The description is built per answer type; this is the cycle-check version, with its evidence field.
\begin{lstlisting}
Submit your final answer, which must be true or false. This ends the task.

answer: boolean
cycle (optional): If your answer is true: the nodes of one cycle in order,
  e.g. [0, 1, 2] for the cycle 0-1-2-0.
\end{lstlisting}

\section*{\code{cannot_answer}}
Offered in every run.
\begin{lstlisting}
Use this when the question cannot be answered with the available tools, or the question
itself is invalid. This ends the task.

reason: string
\end{lstlisting}

\section*{\code{read_result}}
Offered once the first handle exists.
\begin{lstlisting}
Read part of a large stored result by its handle (e.g. "result_1"): up to `limit` items
starting at `offset`.
\end{lstlisting}

\section*{\code{run_python} (mode \code{tools})}
Offered when the sandbox is on; in mode \code{networkx} the sentence ``Only these functions are available: networkx cannot be imported.'' is replaced by one saying that \code{G} (the graph) and \code{nx} (networkx) are also available.
\begin{lstlisting}
Run Python code against G in a sandbox. The graph tools are available as functions with
the same arguments, returning plain values: get_neighbors(4) -> [2, 5], degree(4) -> 2,
has_edge(0, 1) -> True, shortest_path(0, 5) -> [0, 2, 5] (None if there is no path),
graph_info() -> {'nodes': 29, 'edges': 148, 'directed': False} (a dict, like
is_bipartite()). A tool error raises ValueError. Don't redefine these functions. Only
these functions are available: networkx cannot be imported. Variables persist between
calls. Put the value you want back in a variable named `result`; printed output is shown
too (shortened if long). No network or files; each call may run at most 10 seconds.
\end{lstlisting}
